\begin{table*}[tb]
  \centering
  \caption{Feature-space distribution of FoodLogAthl-218 and UEC-Food 256 under three feature extractors. Higher Trace and Davies--Bouldin (DB) indicate larger intra-class spread; lower Silhouette and Calinski--Harabasz (CH) indicate less compact, more overlapping classes. Arrows give the direction that means ``harder''; \textbf{bold} marks the harder dataset in each column.}
  \footnotesize
  \setlength{\tabcolsep}{3pt}
  \begin{tabular}{@{}lcccccccccccc@{}}
      \toprule
        & \multicolumn{3}{c}{\textbf{Trace} ($\uparrow$)}
        & \multicolumn{3}{c}{\textbf{Silhouette} ($\downarrow$)}
        & \multicolumn{3}{c}{\textbf{DB} ($\uparrow$)}
        & \multicolumn{3}{c}{\textbf{CH} ($\downarrow$)} \\
      \cmidrule(lr){2-4} \cmidrule(lr){5-7} \cmidrule(lr){8-10} \cmidrule(l){11-13}
      \textbf{Dataset}
      & CLIP & EVA-CLIP & Swin
      & CLIP & EVA-CLIP & Swin
      & CLIP & EVA-CLIP & Swin
      & CLIP & EVA-CLIP & Swin \\
      \midrule
      UEC-Food 256
        & 71.21            & 45.09            & 264.00
        & 0.0335           & 0.0577           & $-$0.0064
        & 3.64             & 3.02             & 4.11
        & 74.05            & 115.80           & 49.17            \\
      FoodLogAthl-218 (ours)
        & \textbf{81.70}   & \textbf{49.92}   & \textbf{282.57}
        & \textbf{$-$0.0210} & \textbf{0.0036}  & \textbf{$-$0.0288}
        & \textbf{4.52}    & \textbf{3.68}    & \textbf{4.51}
        & \textbf{27.51}   & \textbf{51.51}   & \textbf{21.29}   \\
      \bottomrule
    \end{tabular}
  \label{table:diversity_metrics_comparison}
\end{table*}
